\chapter{Introdução} \label{cap:01}

O acesso rápido a informações confiáveis é importante para o acompanhamento de processos e para a tomada de decisões nas organizações. Entretanto, os dados corporativos costumam estar armazenados em bancos de dados relacionais cuja consulta exige conhecimento de linguagens específicas, como a Structured Query Language (SQL). Em consequência, colaboradores que conhecem o contexto do negócio, mas não dominam SQL, frequentemente dependem de profissionais de Tecnologia da Informação (TI) para obter relatórios e responder a perguntas operacionais. Essa dependência pode aumentar o tempo necessário para acessar informações e direcionar a equipe técnica a solicitações repetitivas de extração de dados.

Os avanços em Inteligência Artificial (IA), especialmente no Processamento de Linguagem Natural (PLN), ampliaram as possibilidades de interação entre pessoas e sistemas computacionais. Nesse contexto, a área denominada \textit{text-to-SQL} investiga métodos capazes de interpretar uma pergunta formulada em linguagem natural e transformá-la em uma consulta SQL compatível com a estrutura de um banco de dados \cite{deng2022recent}. Essa abordagem pode tornar a consulta de dados mais acessível, mas ainda envolve desafios relacionados à compreensão da intenção do usuário, à associação entre os termos da pergunta e o esquema do banco, à correção da consulta produzida e à generalização para diferentes estruturas de dados \cite{yu2018spider}.

Este trabalho tem como tema a aplicação de IA na consulta a bancos de dados e delimita-se ao desenvolvimento de uma aplicação desktop para consultas a bancos de dados PostgreSQL. Por meio da plataforma, o usuário poderá formular, em português e sem escrever SQL, perguntas como ``Quais foram os cinco produtos mais vendidos em julho?''. A aplicação enviará a pergunta e as informações necessárias sobre o esquema do banco a um modelo de IA, receberá uma consulta SQL, verificará sua sintaxe e sua segurança, executará apenas consultas de leitura validadas e explicitamente confirmadas pelo usuário e apresentará o resultado de forma compreensível.

O problema que orienta o projeto pode ser expresso pela seguinte questão: como permitir que usuários sem conhecimento técnico em SQL consultem dados relacionais por meio de linguagem natural, preservando a segurança, a integridade e o controle de acesso às informações? Para responder a essa questão, a solução proposta combinará geração de SQL por IA com uma camada de validação independente da resposta do modelo. Serão aceitas exclusivamente consultas do tipo \texttt{SELECT}, e comandos capazes de alterar estruturas ou registros, como \texttt{INSERT}, \texttt{UPDATE}, \texttt{DELETE}, \texttt{DROP}, \texttt{ALTER} e \texttt{TRUNCATE}, deverão ser bloqueados antes da execução.

O MVP contemplará cadastro e login da conta do produto, configuração de conexões PostgreSQL, escolha entre Ollama local e API externa com chave do próprio usuário, histórico local de conversas e apresentação dos resultados em tabelas. A preparação do contexto combinará a estrutura acessível do banco com regras de negócio confirmadas pelo usuário em uma entrevista inicial por conexão. A interface utilizará React, Vite e TypeScript, distribuídos com Electron, e o \textit{back-end} local será implementado em Python com FastAPI. Gráficos, explicações textuais dos resultados, favoritos, sugestões de perguntas e exportações avançadas não integrarão o compromisso da V1. O sistema será avaliado em cenários controlados e com bancos de dados de teste, sem permitir alterações nos dados consultados.

\section{Justificativa}

Sob a perspectiva teórica, o projeto se justifica pela relevância do problema \textit{text-to-SQL} para as áreas de PLN e bancos de dados. A tradução de uma pergunta para SQL não corresponde apenas à substituição de palavras: o sistema precisa identificar a intenção expressa, relacionar entidades e atributos ao esquema disponível e construir uma consulta sintática e semanticamente adequada. A literatura apresenta avanços importantes nessa área, mas também reconhece desafios de interpretação, geração e avaliação das consultas \cite{deng2022recent}. O conjunto de dados Spider, por exemplo, foi elaborado para avaliar consultas complexas e a generalização entre diferentes domínios e esquemas de bancos de dados, evidenciando a necessidade de soluções que sejam testadas além de exemplos simples ou de uma única base \cite{yu2018spider}.

Em termos acadêmicos, o desenvolvimento da plataforma possibilita integrar conhecimentos de engenharia de software, bancos de dados, segurança da informação, desenvolvimento web e IA em um problema aplicado. Além de construir um protótipo funcional, o trabalho permitirá observar em quais tipos de pergunta a geração automática produz resultados corretos, quais erros ocorrem com maior frequência e em que medida as verificações implementadas impedem consultas inadequadas. Assim, a avaliação do sistema poderá produzir evidências sobre as possibilidades e as limitações do uso de modelos de IA como interface para dados relacionais.

Sob a perspectiva prática e de mercado, a solução busca aproximar os dados dos profissionais que precisam utilizá-los, mas que não possuem formação técnica em SQL. Ao permitir consultas em linguagem cotidiana, a plataforma poderá reduzir o tempo gasto em solicitações simples de relatórios, dar maior autonomia aos usuários e liberar a equipe de TI para atividades de maior complexidade. A apresentação tabular, acompanhada da possibilidade de inspecionar a SQL gerada, poderá facilitar a conferência dos resultados e ampliar o uso das informações disponíveis no processo decisório.

Por fim, a adoção de IA nesse contexto exige cuidados específicos, pois uma consulta incorreta pode apresentar resultados enganosos, consumir recursos excessivos ou afetar informações sensíveis. Por esse motivo, a proposta não tratará a resposta do modelo como uma instrução automaticamente confiável. A restrição a operações de leitura, a validação prévia do SQL, o controle de permissões, o registro do histórico e o uso de bancos de teste constituem partes centrais da solução. Dessa maneira, o projeto procura conciliar facilidade de uso e inovação com mecanismos objetivos de segurança e rastreabilidade.

\section{Objetivos}

\subsection{Objetivo Geral}

Desenvolver e avaliar um protótipo de aplicação desktop que permita a usuários sem domínio de SQL consultar bancos de dados PostgreSQL por meio de perguntas em linguagem natural, utilizando IA para gerar consultas de leitura, validando-as e solicitando confirmação explícita antes da execução e apresentando os resultados de maneira segura e compreensível.

\subsection{Objetivos Específicos}

\begin{itemize}
    \item Revisar conceitos, técnicas e trabalhos relacionados a PLN, modelos de IA e conversão de linguagem natural para SQL;
    \item Definir a arquitetura da plataforma e o fluxo de comunicação entre interface desktop, \textit{back-end}, modelo de IA e banco de dados;
    \item Desenvolver uma interface desktop intuitiva com tecnologias web para o envio de perguntas e a visualização das respostas, considerando usuários sem conhecimento técnico em SQL;
    \item Implementar o \textit{back-end} majoritariamente em Python, adotando FastAPI como camada de API local;
    \item Integrar modelos de IA por Ollama local e por API externa do usuário, capazes de gerar SQL a partir de perguntas em português, do esquema autorizado e do contexto de negócio da conexão;
    \item Implementar uma camada independente de validação sintática e de segurança que autorize somente consultas de leitura e bloqueie comandos de alteração de dados ou de estrutura;
    \item Implementar cadastro e login da conta do produto, respeitar as permissões da role PostgreSQL preexistente e persistir configurações e histórico de conversas localmente;
    \item Apresentar a SQL para inspeção e confirmação explícita do usuário e exibir os resultados em tabelas;
    \item Conduzir uma entrevista inicial para registrar regras de negócio validadas em um arquivo de contexto por conexão e permitir a correção controlada de SQL dentro do limite de três tentativas totais; e
    \item Avaliar o protótipo em bancos de dados de teste, considerando a correção das consultas geradas, a efetividade das regras de segurança, o tempo de execução e a compreensão dos resultados apresentados.
\end{itemize}
